\documentclass[11pt]{article}
\usepackage[margin=1in]{geometry}
\usepackage{amsmath,amssymb,bm}
\usepackage{booktabs}
\usepackage{hyperref}

\newcommand{\vb}[1]{\bm{#1}}
\newcommand{\code}[1]{\texttt{#1}}

\title{Rigid Bodies in MARS:\\Grid Interactions, Equations of Motion, and Energy Bookkeeping}
\date{October 2026}

\begin{document}
\maketitle

\section{Notation and frames}

A rigid body $b$ has lab-frame position $\vb r_b$, orientation
$\mathsf R_b\in SO(3)$ (columns are the body axes in the lab frame), linear
momentum $\vb p_b$ (lab frame), and angular momentum $\vb L_b$ stored in the
\emph{body} frame. Its type supplies a mass $m_b$ and principal moments
$\vb I_b=(I_x,I_y,I_z)$; the body frame is assumed principal.

A grid $G$ attached to a body has an origin $\vb o_G$ and a basis matrix
$\mathsf B_G$ (columns are the voxel step vectors), both in the body frame. The
voxel with integer index $\vb n=(n_x,n_y,n_z)$ sits at
\begin{equation}
  \vb x_{\vb n} \;=\; \mathsf R_b\left(\mathsf B_G\,\vb n + \vb o_G\right) + \vb r_b
  \;=\; \mathsf A_G\,\vb n + \vb O_G,
  \qquad
  \mathsf A_G \equiv \mathsf R_b\mathsf B_G,\quad
  \vb O_G \equiv \mathsf R_b\vb o_G + \vb r_b ,
  \label{eq:voxel}
\end{equation}
where $\mathsf A_G$ and $\vb O_G$ are the grid's lab-frame basis and origin.
A lab point $\vb x$ maps to continuous index coordinates
$\vb s = \mathsf A_G^{-1}(\vb x-\vb O_G)$. A grid owned by no body (an external
PMF) uses $\mathsf R=\mathsf 1$, $\vb r=\vb 0$.

Grid values are interpolated in index space (linear or cubic), giving
$\tilde u(\vb s)$ and $\nabla_{\vb s}\tilde u$. By the chain rule the
lab-frame gradient is
\begin{equation}
  \nabla_{\vb x} u = \mathsf A_G^{-\mathsf T}\,\nabla_{\vb s}\tilde u .
  \label{eq:chain}
\end{equation}

Every grid enters with a scale factor $\lambda$ (\code{GridTerm::scale},
config \code{gridFileScale}).

\paragraph{Reference point.} All torques on body $b$ are taken about
$\vb r_b$, because the integrators (Section~\ref{sec:eom}) apply
$\mathsf R_b^{\mathsf T}\vb\tau_b$ to the body-frame angular momentum. If a
torque $\vb\tau(\vb c)$ is computed about a point $\vb c$ for forces with sum
$\vb F$, then
\begin{equation}
  \vb\tau(\vb r_b) \;=\; \vb\tau(\vb c) + (\vb c-\vb r_b)\times\vb F .
  \label{eq:shift}
\end{equation}

\section{Grid--grid interaction between two bodies}
\label{sec:gridgrid}

\subsection{Energy}

Body $A$ carries a \emph{density} grid $\rho$ and body $B$ a \emph{potential}
grid $u$. Grid pairs are matched by key name: a density grid of type $t_i$
pairs with a potential grid of type $t_j$ with the same key, for $t_i\le t_j$
(the legacy one-direction rule, kept on purpose; \code{RigidBodyForcePairList}).
The interaction energy is the overlap integral, discretized over $\rho$'s
voxels:
\begin{equation}
  E_{AB} \;=\; \lambda\int \rho(\vb x)\,u(\vb x)\,d^3x
  \;\approx\; \lambda\sum_{\vb n}\rho_{\vb n}\,
  \tilde u\!\left(\vb s_{\vb n}\right),
  \qquad
  \vb s_{\vb n} = \mathsf A_u^{-1}\!\left(\vb x_{\vb n}-\vb O_u\right),
  \label{eq:Egg}
\end{equation}
with $\vb x_{\vb n}$ from Eq.~\eqref{eq:voxel} for grid $\rho$ on body $A$. In
the kernel the index map is
$\vb s_{\vb n} = \mathsf A_u^{-1}\left(\mathsf A_\rho\vb n + \vb\Delta\right)$
with $\vb\Delta=\vb O_\rho-\vb O_u$ (\code{origin\_offset}).

\emph{Note.} v1 stored $\sum_{\vb n}\rho_{\vb n}$ as the energy, which does
not depend on the bodies' relative pose. MARS uses Eq.~\eqref{eq:Egg}, the only
energy whose negative gradient is the force both codes compute.

\subsection{\texorpdfstring{Force and torque on $A$}{Force and torque on A}}

Each voxel is a point load. Its force is
\begin{equation}
  \vb f_{\vb n} \;=\; -\lambda\,\rho_{\vb n}\,\nabla_{\vb x}u(\vb x_{\vb n})
  \;=\; -\lambda\,\rho_{\vb n}\,\mathsf A_u^{-\mathsf T}\nabla_{\vb s}\tilde u(\vb s_{\vb n}),
\end{equation}
and
\begin{equation}
  \begin{aligned}
  \vb F_A&=\sum_{\vb n}\vb f_{\vb n},\qquad
  \vb\tau_A(\vb O_\rho)=\sum_{\vb n}\left(\mathsf A_\rho\vb n\right)\times\vb f_{\vb n},\\
  \vb\tau_A(\vb r_A)&=\vb\tau_A(\vb O_\rho)+(\vb O_\rho-\vb r_A)\times\vb F_A .
  \end{aligned}
  \label{eq:FA}
\end{equation}
The kernel accumulates $\vb\tau_A$ about $\vb O_\rho$ (the lever arm
$\mathsf A_\rho\vb n$ is the voxel's offset from the density grid's origin),
then applies the shift~\eqref{eq:shift} once per block with
$\vb O_\rho-\vb r_A=\mathsf R_A\vb o_\rho$ (\code{rho\_shift}).

\subsection{\texorpdfstring{Reaction on $B$}{Reaction on B}}

$E_{AB}$ is unchanged when both bodies are translated or rotated together, so
the discretized sum conserves total linear and angular momentum exactly. The
force and torque on $B$ therefore follow in closed form, with no second grid
pass:
\begin{equation}
  \vb F_B=-\vb F_A,\qquad
  \vb\tau_B(\vb r_B)=-\vb\tau_A(\vb r_A)+(\vb r_B-\vb r_A)\times\vb F_A .
  \label{eq:reaction}
\end{equation}
The kernel uses the equivalent form
$\vb\tau_B(\vb O_u)=-\left[\vb\tau_A(\vb O_\rho)+\vb\Delta\times\vb F_A\right]$,
shifted to $\vb r_B$ with $\vb O_u-\vb r_B=\mathsf R_B\vb o_u$
(\code{u\_shift}). Expanding $\vb\Delta$ shows it equals
Eq.~\eqref{eq:reaction}. It needs no body positions, and it uses the same
(unwrapped) $\vb\Delta$ as the force.

\subsection{External PMF grids}

A density grid can also be matched against one of its own type's PMF grids.
The PMF is fixed in the lab frame ($\mathsf A_u=\mathsf B_u$,
$\vb O_u=\vb o_u$). There is no second body, so Eq.~\eqref{eq:reaction} does
not apply and $A$ keeps all of $E$.

\subsection{Update period}

A pair may be evaluated only every $P$ steps (\code{rigidBodyGridGridPeriod}).
For the time-averaged impulse to be the same as evaluating every step, the
force and torque applied on an evaluation step should be multiplied by $P$, as
v1 does.

\subsection{Kernels}

\code{RBGridCullKernel} keeps (instance pair, grid pair) candidates within the
cutoff, or always keeps them for a PMF, and precomputes
$\mathsf A_\rho$, $\mathsf A_u^{-1}$ and $\vb\Delta$.
\code{RBGridPrefixSumKernel} assigns block ranges.
\code{RBGridBatchedForceKernel} splits each item's $\rho$ voxels over its
blocks, evaluates \code{grid\_grid\_voxel\_force\_torque}, block-reduces
$(\vb F,E)$ and $\vb\tau$, and atomically adds the results into the bodies
(\code{RigidBodyGridBatch.h}, \code{GridGridKernels.h}).

\section{Particles in a body's potential grid}
\label{sec:pgrid}

\subsection{Selection}

A particle type samples body potential grid $u$ only if its \code{particle}
block lists the grid's key with \code{rigidBodyPotential}. At dispatch time
one candidate is made per (body instance, potential grid) for which at least
one particle type matches, along with a per-type flag row; the kernel skips
particles whose type flag is 0. This is v1's \code{partRigidBodyGrid} rule.
A type mask is used instead of per-grid particle-index lists because it stays
valid when particles are reordered.

\subsection{Force, torque, energy}

For a selected particle $k$ at $\vb x_k$,
with $\vb s_k=\mathsf A_u^{-1}(\vb x_k-\vb O_u)$:
\begin{align}
  U_k &= \lambda\,\tilde u(\vb s_k), &
  \vb f_k &= -\lambda\,\mathsf A_u^{-\mathsf T}\nabla_{\vb s}\tilde u(\vb s_k),\\
  \vb F_b &= -\sum_k\vb f_k, &
  \vb\tau_b(\vb O_u) &= \sum_k(\vb x_k-\vb O_u)\times(-\vb f_k),
\end{align}
and $\vb\tau_b(\vb r_b)=\vb\tau_b(\vb O_u)+(\vb O_u-\vb r_b)\times\vb F_b$ by
Eq.~\eqref{eq:shift}, with $\vb O_u-\vb r_b=\mathsf R_b\vb o_u$
(\code{grid\_shift}). Each particle receives $\vb f_k$ by atomic scatter,
because several candidates can act on the same particle. The body's totals are
block-reduced, then added atomically
(\code{RBParticleGridForceKernel}, \code{RigidBodyParticleGridBatch.h}).

\section{Attached particles}
\label{sec:attached}

An attached particle $k$ is an ordinary particle held at a fixed body-frame
offset $\vb d_k$:
\begin{equation}
  \vb x_k=\mathsf R_b\vb d_k+\vb r_b .
\end{equation}
Its position is rewritten from the body each step
(\code{RBSyncAttachedPositionsKernel}) and is not integrated. Every force path
(pairwise, bonded, PMF, body grids) acts on it as on any other particle. The
total force on it is then transferred to the body
(\code{RBReduceAttachedForcesKernel}):
\begin{equation}
  \vb F_b \mathrel{+}= \sum_k\vb f_k,\qquad
  \vb\tau_b(\vb r_b) \mathrel{+}= \sum_k(\mathsf R_b\vb d_k)\times\vb f_k .
\end{equation}
The lever arm is measured from $\vb r_b$, so $\vb r_b$ must be the point the
attached template is built about (\code{referencePoint}). The particle's own
force is left in place; nothing reads it afterwards, and the force output then
still shows what acted on the particle.

\section{Equations of motion}
\label{sec:eom}

Let $\vb F_b$ and $\vb\tau_b$ (about $\vb r_b$, lab frame) be the totals from
Sections~\ref{sec:gridgrid}--\ref{sec:attached}, plus any host-set external
load. The body-frame torque is $\mathsf R_b^{\mathsf T}\vb\tau_b$.

\paragraph{Langevin / DLM (\code{RBDLM.h}).}
Langevin forces are added first (\code{RBLangevinForceKernel}). In the body
frame, per axis $\alpha$,
\begin{equation}
  f_\alpha = \sqrt{\frac{2k_BT\,m\,\gamma_\alpha}{\Delta t}}\,\xi_\alpha
  - \gamma_\alpha\left(\mathsf R^{\mathsf T}\vb p\right)_\alpha ,
  \qquad
  \tau_\alpha = \sqrt{\frac{2k_BT\,I_\alpha\,\Gamma_\alpha}{\Delta t}}\,\eta_\alpha
  - \Gamma_\alpha L_\alpha ,
\end{equation}
with $\xi,\eta\sim\mathcal N(0,1)$ and translational and rotational damping
$\gamma,\Gamma$. These are rotated to the lab frame and added to
$\vb F_b,\vb\tau_b$. The dissipative term carries legacy unit factors
(\code{langevin\_damping\_unit}, \code{langevin\_damp\_scale}). The step is
split DLM (Dullweber--Leimkuhler--McLachlan): half kick
$\vb p\mathrel{+}=\tfrac{\Delta t}{2}\vb F$,
$\vb L\mathrel{+}=\tfrac{\Delta t}{2}\mathsf R^{\mathsf T}\vb\tau$; drift
$\vb r\mathrel{+}=\Delta t\,\vb p/m$; free rotation by the symmetric splitting
$x(\tfrac{\Delta t}{2})\,y(\tfrac{\Delta t}{2})\,z(\Delta t)\,y(\tfrac{\Delta t}{2})\,x(\tfrac{\Delta t}{2})$,
each factor being the exact rotation about one body axis with rate
$L_\alpha/I_\alpha$; then the second half kick after the next force
evaluation.

\paragraph{Brownian (\code{RBBD.h}).}
Overdamped, with no momenta. With mobilities $\mu_\alpha=1/(\gamma_\alpha m)$
and $\mu^{r}_\alpha=1/(\Gamma_\alpha I_\alpha)$ in the body frame,
\begin{equation}
  \begin{aligned}
  \delta\vb r_{\mathrm{body}} &= \mu\odot\mathsf R^{\mathsf T}\vb F\,\Delta t
  + \sqrt{2k_BT\mu\,\Delta t}\odot\vb\xi,\\
  \delta\vb\theta &= \mu^{r}\odot\mathsf R^{\mathsf T}\vb\tau\,\Delta t
  + \sqrt{2k_BT\mu^{r}\Delta t}\odot\vb\eta ,
  \end{aligned}
\end{equation}
then $\vb r\mathrel{+}=\mathsf R\,\delta\vb r_{\mathrm{body}}$ and
$\mathsf R\leftarrow\mathsf R\,\mathcal R(\delta\vb\theta)$, where
$\mathcal R$ is the symmetric $z\,y\,x\,y\,z$ splitting.

\paragraph{Units.} Lengths are \AA, time ns, mass amu, energy kcal/mol. Force
times time becomes momentum through \code{impulse\_to\_momentum}
($=4.1868\times10^4$), and $\vb p/m$ becomes \AA/ns through
\code{velocity\_scale} ($=10^4$). The stored momentum is therefore in units
where $\vb p/m$ is measured in $10^4$\,\AA/ns $=1$\,km/s.

\section{Energy}
\label{sec:energy}

\subsection{Potential energy: shares}

Every interaction energy is divided once among its partners, so the shares add
up to the total with nothing counted twice. Each particle stores its share in
the fourth component of its force (\code{ForceEnergy.t}); each body stores its
share in \code{rb.force[b].t}.

\begin{center}
\begin{tabular}{@{}lll@{}}
\toprule
Interaction & Particle share & Body share\\
\midrule
Nonbonded pair, bond & $\tfrac12$ each & ---\\
Angle / dihedral & $\tfrac13$ / $\tfrac14$ each & ---\\
Restraint, particle PMF & all & ---\\
Particle in body potential grid & $\tfrac12$ & $\tfrac12$\\
Body--body grid pair & --- & $\tfrac12$ each\\
Body density in external PMF & --- & all\\
\bottomrule
\end{tabular}
\end{center}

The split is by partner count, not by mass. Potential energy belongs to the
interaction rather than to either partner, and mass only determines how
\emph{kinetic} energy is shared.

Writing $U^{\mathrm{grid}}_b=$ \code{rb.force[b].t}, the system potential
energy is
\begin{equation}
  U_{\mathrm{total}}=\sum_{k\in\text{particles}}U_k+\sum_b U^{\mathrm{grid}}_b .
\end{equation}
A body whose type has no grids has $U^{\mathrm{grid}}_b=0$.

Attached particles are ordinary particles, so their energy is already in
$\sum_k U_k$. It is also reported for each body as
\begin{equation}
  U^{\mathrm{att}}=\sum_{k\in\text{attached}}U_k ,
\end{equation}
because for a body that is just a collection of particles this, not
$U^{\mathrm{grid}}$, is the body's interaction energy. Which term is
meaningful depends on the model, so both are reported and neither is assumed.

\subsection{Kinetic energy}

\begin{equation}
  K_t=\frac{\vb p^2}{2m},\qquad
  K_r=\sum_{\alpha}\frac{L_\alpha^2}{2I_\alpha}\quad(\text{body frame}).
\end{equation}
In stored units $\vb p^2/m$ is in amu\,(km/s)$^2=1$\,kJ/mol, so
\begin{equation}
  K\,[k_BT]=\frac{K_{\mathrm{raw}}}{4.1868\;k_BT[\mathrm{kcal/mol}]},
  \qquad 4.1868=\code{SQRT\_CAL\_TO\_JOULE}^2 ,
\end{equation}
which reproduces v1's factor $2.388458509\times10^{-1}=1/4.1868$ and matches
the particle kinetic energy.

\subsection{Reduction and output}

\code{RBEnergyReduceKernel} (\code{RBOperation/RBEnergyKernels.h}) has thread
$i$ handle body $i$ and attached particle $i$. It block-reduces one
\code{RB\_Energy} $=(K_t,K_r,U^{\mathrm{grid}},U^{\mathrm{att}})$ and
atomically adds it into a single device value, so one 16-byte copy reaches the
host on each energy-output step. \code{rb\_energy.dat}:
\begin{verbatim}
Kinetic Energy <Kt+Kr> (kT)
  Translational <Kt>  Rotational <Kr> (kT)
Potential Energy <Ugrid+Uatt> (kcal/mol)
  Grid <Ugrid>  Attached <Uatt> (kcal/mol)
\end{verbatim}
The first and third lines keep the v1 format.

\section{Implementation status}
\label{sec:status}

The code departs from the equations above in these places:
\begin{enumerate}
  \item \textbf{Update period.} Forces and torques from pairs with $P>1$ are
    not multiplied by $P$.
  \item \textbf{SMD scale slope.} \code{gridFileScaleSlope} is not applied to
    body potential grids.
  \item \textbf{Periodicity.} Grid--grid offsets are not wrapped; periodic
    body--body pairs are future work.
\end{enumerate}

\end{document}
